
\section{\textbf{\texttt{get\_view()}} ORM Method}

\begin{frame}[fragile]{\textbf{\texttt{get\_view()}} Method}
	\begin{itemize}
		\item The \textbf{\texttt{get\_view()}} method returns the architecture and metadata of a view for a model (form, list/tree, kanban, search, \ldots).
		\item The web client uses it to know how to render a screen.
		\item In Odoo 17, \texttt{get\_view()} replaces older multi-view helpers such as \texttt{fields\_view\_get()} in many flows.
\begin{block}{Method Syntax}
\begin{lstlisting}
@api.model
def get_view(self, view_id=None, view_type='form', **options):
	return super().get_view(view_id=view_id, view_type=view_type, **options)
\end{lstlisting}
\end{block}
	\end{itemize}
\end{frame}

\begin{frame}[fragile]{Breaking Down the Method Signature}
	\begin{block}{}
		\begin{itemize}
			\item \textbf{\texttt{view\_id}}: Optional specific view XML ID / database ID.
			\item \textbf{\texttt{view\_type}}: Type of view to load: \texttt{'form'}, \texttt{'list'}, \texttt{'kanban'}, \texttt{'search'}, \ldots
			\item \textbf{\texttt{options}}: Extra options used by the client / ORM.
\begin{lstlisting}
view_type = 'form'
view_id = None  # let Odoo choose the default form view
\end{lstlisting}
		\end{itemize}
	\end{block}
\end{frame}

\begin{frame}[fragile]{What Does \texttt{get\_view()} Return?}
	\begin{block}{}
		\begin{itemize}
			\item It returns a dictionary of view information.
\begin{lstlisting}
info = self.env['student.student'].get_view(view_type='form')
print(info.keys())
\end{lstlisting}
			Common keys include:
			\begin{itemize}
				\item \texttt{id} -- view record ID
				\item \texttt{name} -- view name
				\item \texttt{type} -- view type
				\item \texttt{arch} -- the XML architecture
				\item \texttt{models} -- field/model metadata needed by the view
			\end{itemize}
		\end{itemize}
	\end{block}
\end{frame}

\begin{frame}[fragile]{Practical Example}
\begin{lstlisting}
form_view = self.env['student.student'].get_view(view_type='form')
list_view = self.env['student.student'].get_view(view_type='list')

print(form_view['type'])   # form
print(list_view['type'])   # list
print(form_view['arch'])   # XML architecture
\end{lstlisting}
	\begin{itemize}
		\item Developers rarely call this in business methods.
		\item It is essential for UI, studio-like tools, and technical debugging.
	\end{itemize}
\end{frame}

\begin{frame}[fragile]{When Is \texttt{get\_view()} Useful?}
	\begin{itemize}
		\item Understanding which view Odoo will open for a model.
		\item Building custom clients that render Odoo views.
		\item Debugging inherited / extended view architecture.
		\item Inspecting the final combined XML after view inheritance.
\begin{lstlisting}
info = self.env['student.student'].get_view(view_type='search')
\end{lstlisting}
	\end{itemize}
\end{frame}

\begin{frame}[fragile]{Method Call}
\begin{block}{Method Call}
\begin{lstlisting}
view_info = self.env['student.student'].get_view(
	view_id=None,
	view_type='form',
)
\end{lstlisting}
	\begin{itemize}
		\item Here, \textbf{\texttt{.get\_view(...)}} is the \textbf{method call}.
		\item Return type reminder: \textbf{dict} of view metadata.
	\end{itemize}
\end{block}
\end{frame}
